\chapter{uv のインストール}

\section{インストールコマンド}

WSL のターミナルを開いて以下を実行します。
\textbf{これは最初の1回だけ}でよいです。

\begin{wslbox}[uv のインストール]
\begin{lstlisting}[style=bash]
$ curl -LsSf https://astral.sh/uv/install.sh | sh
\end{lstlisting}
\end{wslbox}

実行すると以下のような出力が出ます。

\begin{wslbox}[インストール時の出力例]
\begin{lstlisting}[style=bash]
downloading uv 0.6.x x86_64-unknown-linux-musl
installing to /home/username/.local/bin
  uv
  uvx
everything's installed!

To add $HOME/.local/bin to your PATH, either
restart your shell or run:

    source $HOME/.bashrc
\end{lstlisting}
\end{wslbox}

\begin{notebox}[\cmd{curl ... | sh} について]
インターネットからスクリプトをダウンロードして直接実行するコマンドです。
Astral 社の公式インストーラーなので安全です。
\end{notebox}

\section{PATH の反映}

インストール直後はコマンドがまだ認識されません。
\textbf{ターミナルを一度閉じて開き直してください。}

あるいは、ターミナルを開き直さずに即座に反映させる場合は：

\begin{wslbox}[PATH を今すぐ反映する]
\begin{lstlisting}[style=bash]
$ source ~/.bashrc
\end{lstlisting}
\end{wslbox}

\section{インストールの確認}

以下のコマンドでバージョンを確認します。

\begin{wslbox}[バージョン確認]
\begin{lstlisting}[style=bash]
$ uv --version
uv 0.6.x (linux x86_64 ...)
\end{lstlisting}
\end{wslbox}

このように表示されれば OK です。

\begin{wslbox}[失敗した場合の出力]
\begin{lstlisting}[style=bash]
$ uv --version
-bash: uv: command not found
\end{lstlisting}
\end{wslbox}

\cmd{command not found} と表示された場合は、
ターミナルの再起動か \cmd{source ~/.bashrc} を試してください。
それでも解決しない場合は第~\ref{chap:errors}章を参照してください。

\section{uv のアップデート（将来の参考用）}

uv 自体のアップデートは以下のコマンドで行えます。
セットアップ時に実行する必要はありません。

\begin{wslbox}[uv のアップデート]
\begin{lstlisting}[style=bash]
$ uv self update
\end{lstlisting}
\end{wslbox}
